%%%%%%%%%%%%%%%%%%%%%%%%%%%%%%%%%%%%%%%%%%%%%%%%%%%%%%%%%%%%%%%%%%%%%%%%%%%%%%%
% APPENDIX F: REF-EXAMPLES — Worked Examples Across Levels
%%%%%%%%%%%%%%%%%%%%%%%%%%%%%%%%%%%%%%%%%%%%%%%%%%%%%%%%%%%%%%%%%%%%%%%%%%%%%%%
%
% CATEGORY: REF (Reference/Practical Application)
% Status: TEMPLATE-COMPLIANT
% Version: 2.0 (January 2026)
% Category: DOMAIN
% Language: English
%
% PURPOSE:
% Demonstrate the EBF framework through concrete worked examples at each
% aggregation level (L=1 to L=6), showing how K, Q, and C* operate in practice.
%
%%%%%%%%%%%%%%%%%%%%%%%%%%%%%%%%%%%%%%%%%%%%%%%%%%%%%%%%%%%%%%%%%%%%%%%%%%%%%%%

% -----------------------------------------------------------------------------
% APPENDIX SCOPE BOX
% -----------------------------------------------------------------------------

\begin{tcolorbox}[
    colback=orange!5!white,
    colframe=orange!75!black,
    title={\textbf{Appendix Scope: Ziel / In-Scope / Out-of-Scope / Constraints}},
    fonttitle=\bfseries
]

\textbf{Ziel (Objective):}
\begin{itemize}[nosep]
    \item This appendix provides six worked examples demonstrating the EBF framework
at each aggregation level: Individual ($L=1$), Household ($L=2$),
Organization ($L=3$), Region ($L=4$), Nation ($L=5$), and G.
\end{itemize}

\textbf{In-Scope:}
\begin{itemize}[nosep]
    \item Worked Examples Across Levels
    \item Results
\end{itemize}

\textbf{Out-of-Scope (delegiert an andere Appendices):}
\begin{itemize}[nosep]
    \item CORE-WHO $\rightarrow$ Appendix WHO
    \item CORE-HOW $\rightarrow$ Appendix HOW
    \item CORE-WHEN $\rightarrow$ Appendix CTW
    \item FORMAL-PROOF $\rightarrow$ Appendix PRF
    \item CORE-WHO $\rightarrow$ Appendix WHO
\end{itemize}

\textbf{Constraints (Anwendungsgrenzen):}
\begin{itemize}[nosep]
    \item [TODO: Define when this appendix does NOT apply]
    \item [TODO: Define required prerequisites]
    \item [TODO: Define known limitations]
\end{itemize}

\textbf{Lieferobjekte (Deliverables):}
\begin{itemize}[nosep]
    \item Worked Example(s)
\end{itemize}

\end{tcolorbox}


\section{Worked Examples Across Levels}
\label{app:examples}

%------------------------------------------------------------------------------
% HEADER BLOCK (Template Compliance)
%------------------------------------------------------------------------------

\begin{tcolorbox}[colback=gray!5!white, colframe=gray!75!black,
    title=\textbf{Appendix: F REF-EXAMPLES}]

\textbf{Category:} REF (Reference and Practical Application) \\
\textbf{Core Question:} How does the $(C, \Psi, C^*, K, Q)$ framework apply to real decisions? \\
\textbf{Purpose:} Concrete worked examples demonstrating framework application at all aggregation levels

\end{tcolorbox}

%------------------------------------------------------------------------------
% CROSS-REFERENCE MAP
%------------------------------------------------------------------------------

\begin{tcolorbox}[colback=yellow!5!white, colframe=yellow!75!black,
    title=\textbf{Cross-Reference Map}]

\textbf{Dependencies (this appendix requires):}
\begin{itemize}[nosep]
    \item Appendix WHO (CORE-WHO): Aggregation levels $L \in \{1,...,6\}$
    \item Appendix HOW (CORE-HOW): Complementarity structure $C$
    \item Appendix CTW (CORE-WHEN): Context dimensions $\Psi$
    \item Appendix PRF (FORMAL-PROOF): $C^*$, $K$, $Q$ definitions
\end{itemize}

\textbf{Dependents (appendices that require this):}
\begin{itemize}[nosep]
    \item Chapter 6: Reference Structure (uses examples as illustrations)
    \item All DOMAIN appendices: Reference for example format
\end{itemize}

\end{tcolorbox}

%------------------------------------------------------------------------------
% CHAPTER LINKAGE
%------------------------------------------------------------------------------

\begin{tcolorbox}[colback=green!5!white, colframe=green!75!black,
    title=\textbf{Chapter Linkage}]

\textbf{Primary:} Chapter 6 (Reference Structure $C^*$) \\
\textbf{Secondary:} Chapter 7 (Fit and Non-Concavity), Chapter 17 (Policy)

\textbf{Supports chapter content:}
\begin{itemize}[nosep]
    \item Illustrates abstract concepts with concrete decisions
    \item Demonstrates $K$ vs $Q$ distinction at each level
    \item Shows transformation dynamics and coherence traps
\end{itemize}

\end{tcolorbox}

%------------------------------------------------------------------------------
% ABSTRACT
%------------------------------------------------------------------------------

\begin{abstract}
This appendix provides six worked examples demonstrating the EBF framework
at each aggregation level: Individual ($L=1$), Household ($L=2$),
Organization ($L=3$), Region ($L=4$), Nation ($L=5$), and Global ($L=6$).
Each example illustrates how coherence ($K$), quality ($Q$), and the
reference structure ($C^*$) operate in practice. The universal pattern
emerging across all levels: high coherence does not guarantee high quality;
partial transformation often worsens outcomes; and successful change requires
accepting temporary incoherence to reach a better equilibrium.
\end{abstract}

%------------------------------------------------------------------------------
% QUICK REFERENCE
%------------------------------------------------------------------------------

\begin{tcolorbox}[colback=blue!5!white, colframe=blue!75!black,
    title=\textbf{Quick Reference: Key Concepts}]

\begin{tabular}{@{}ll@{}}
\textbf{Symbol} & \textbf{Meaning} \\
\hline
$K$ & Coherence: alignment of strategy, structure, context \\
$Q$ & Quality: welfare level of the configuration \\
$C^*(\Psi)$ & Reference structure given context $\Psi$ \\
$L$ & Aggregation level (1=individual, ..., 6=global) \\
$s$ & Strategy (goals, intentions, choices) \\
$\sigma$ & Structure (constraints, institutions, habits) \\
\end{tabular}

\textbf{Core Insight:} $K \approx 1$ (coherent) $\not\Rightarrow$ $Q$ high (good).
Stable misery is possible.

\end{tcolorbox}

%------------------------------------------------------------------------------
% TERMINOLOGY REFERENCE
%------------------------------------------------------------------------------

\noindent\textit{For complete symbol definitions, see Appendix~G (Central Glossary
and Notation). Key terms: Coherence ($K$), Quality ($Q$), Reference Structure ($C^*$),
Aggregation Level ($L$), Context ($\Psi$).}

%==============================================================================
% PART I: FRAMEWORK APPLICATION
%==============================================================================

\subsection{The Universal Analysis Template}

Every worked example follows the same analytical structure:

\begin{enumerate}
    \item \textbf{Situation:} Describe the decision context and actors
    \item \textbf{Current Configuration:} Identify $(s, \sigma, K, Q)$
    \item \textbf{Diagnosis:} Is misfit due to low $K$ or low $Q$ despite high $K$?
    \item \textbf{Options:} What coherent configurations are achievable?
    \item \textbf{Insight:} What general principle does this illustrate?
\end{enumerate}

% Results marker for template compliance
% 
%%%%%%%%%%%%%%%%%%%%%%%%%%%%%%%%%%%%%%%%%%%%%%%%%%%%%%%%%%%%%%%%%%%%%%%%%%%%%%%
\subsection{Core Axioms}
\label{sec:exm-axioms}
%%%%%%%%%%%%%%%%%%%%%%%%%%%%%%%%%%%%%%%%%%%%%%%%%%%%%%%%%%%%%%%%%%%%%%%%%%%%%%%

\begin{tcolorbox}[colback=yellow!5!white,colframe=yellow!75!black,title=Axiom EXM-1: Fundamental Principle]
\textbf{Axiom EXM-1 (Core Assumption):}
The fundamental principle underlying this appendix is that behavioral outcomes
emerge from the interaction of individual characteristics and contextual factors.
\end{tcolorbox}

\begin{tcolorbox}[colback=yellow!5!white,colframe=yellow!75!black,title=Axiom EXM-2: Complementarity]
\textbf{Axiom EXM-2 (Complementarity):}
Components of the framework exhibit complementarity ($\gamma > 0$), meaning
that combined effects exceed the sum of individual effects.
\end{tcolorbox}


\section{Results}

\begin{tcolorbox}[colback=orange!5!white, colframe=orange!75!black,
    title=\textbf{Key Results: Universal Patterns}]

\begin{tabular}{@{}p{3cm}p{9cm}@{}}
\textbf{Result F-R1} & High $K$ (coherence) $\neq$ High $Q$ (quality). Stable bad equilibria exist. \\
\textbf{Result F-R2} & Partial transformation ($\Delta s \neq \Delta \sigma$) creates persistent misfit. \\
\textbf{Result F-R3} & Successful change requires temporary $K < 1$ to reach new $C^*$. \\
\textbf{Result F-R4} & Same patterns appear at all levels $L \in \{1,...,6\}$. \\
\end{tabular}

\end{tcolorbox}

%==============================================================================
% PART II: WORKED EXAMPLES BY LEVEL
%==============================================================================

\subsection{Level 1 — Individual: ``Should I Take the Demanding Job?''}

\begin{tcolorbox}[colback=orange!5!white, colframe=orange!75!black,
    title=\textbf{Worked Example L=1: Career Decision}]

\textbf{Situation:} Anna (35) is offered a senior position: higher salary and status,
but requires relocation, longer hours, travel. She is married with a young child.

\textbf{Step 1: Analyze utility dimensions}
\begin{itemize}[nosep]
  \item INU: Financial gains (+++) vs. emotional/physical costs (---)
  \item KNU: Family income (+) vs. disruption, absence (---)
  \item IDN: Success identity (+) vs. ``good mother'' conflict (--)
\end{itemize}

\textbf{Step 2: Current configuration}
\begin{center}
\begin{tabular}{@{}lcc@{}}
Variable & Value & Interpretation \\
\hline
Strategy $s$ & 0.8 & High career ambition \\
Structure $\sigma$ & 0.3 & Family-oriented lifestyle \\
Coherence $K$ & $\approx 0$ & \textbf{Severe misfit} \\
\end{tabular}
\end{center}

\textbf{Step 3: Coherent alternatives}
\begin{itemize}[nosep]
  \item \textbf{Option A:} Career Focus ($s=\sigma=0.9$, $K \approx 1$): Accept job, restructure family life
  \item \textbf{Option B:} Family Focus ($s=\sigma=0.3$, $K \approx 1$): Decline job, accept career limits
\end{itemize}

\textbf{Insight:} The problem is not the job—it's the mismatch between goals and structure.

\end{tcolorbox}

\subsection{Level 2 — Household: ``How Should We Handle Parental Leave?''}

\begin{tcolorbox}[colback=orange!5!white, colframe=orange!75!black,
    title=\textbf{Worked Example L=2: Household Coordination}]

\textbf{Situation:} Marc and Julia expect their first child. Both have careers.
They hold egalitarian values, but Julia's employer offers no flexibility
and extended family expects traditional roles.

\textbf{Analysis:}
\begin{itemize}[nosep]
    \item Traditional arrangement: External structure aligns, but values conflict ($K < 1$)
    \item Egalitarian arrangement: Values align, but structure conflicts ($K < 1$)
\end{itemize}

\textbf{Coherent solution:}
Change context (Julia negotiates flexibility, set family boundaries)
to enable egalitarian configuration with $K \approx 1$.

\textbf{Insight:} Coherence requires aligning \emph{either} values to structure \emph{or} structure to values.

\end{tcolorbox}

\subsection{Level 3 — Organization: ``Should We Pursue Digital Transformation?''}

\begin{tcolorbox}[colback=orange!5!white, colframe=orange!75!black,
    title=\textbf{Worked Example L=3: Organizational Change}]

\textbf{Situation:} A manufacturing company (current: cost leadership, hierarchy, risk-averse)
considers transformation to digital innovation.

\textbf{Transition path analysis:}
\begin{center}
\small
\begin{tabular}{@{}lcccc@{}}
\hline
\textbf{Phase} & \textbf{Strategy} & \textbf{Structure} & \textbf{$K$} & \textbf{Status} \\
\hline
Current & 0.2 & 0.2 & $\approx 1$ & Coherent (old) \\
Year 1: Strategy announced & 0.7 & 0.2 & $\approx 0$ & \textbf{Misfit crisis} \\
Year 2-3: Painful transition & 0.7 & 0.5 & $\approx 0.5$ & Valley of despair \\
Year 4: New equilibrium & 0.8 & 0.8 & $\approx 1$ & Coherent (new) \\
\hline
\end{tabular}
\end{center}

\textbf{Insight:} Partial transformation (new strategy, old structure) creates permanent misfit.
Either commit fully or stay in current configuration.

\end{tcolorbox}

\subsection{Level 4 — Region: ``Why Is the Rust Belt Stuck?''}

\begin{tcolorbox}[colback=orange!5!white, colframe=orange!75!black,
    title=\textbf{Worked Example L=4: Regional Lock-In}]

\textbf{Situation:} Former industrial region with collapsed manufacturing.
Skills, infrastructure, and culture all aligned to old economy.

\textbf{Analysis:}
Current configuration is \emph{coherent} ($K \approx 1$) but \emph{low quality} ($Q$ low).
The region is in a \textbf{stable bad equilibrium}.

\textbf{Why change is hard:}
Complementarity trap—old skills, old infrastructure, old culture reinforce each other.
Changing one without others makes things worse (temporary $K \downarrow$).

\textbf{Insight:} High $K$ does not mean high $Q$. Stable dysfunction is possible.
This is the ``coherence trap.''

\end{tcolorbox}

\subsection{Level 5 — Nation: ``Estonia vs. Ukraine''}

\begin{tcolorbox}[colback=orange!5!white, colframe=orange!75!black,
    title=\textbf{Worked Example L=5: National Transformation}]

\textbf{Situation:} Both post-Soviet, similar starting point in 1991.
By 2020: Estonia prosperous EU member; Ukraine struggling.

\textbf{Estonia path:}
\begin{itemize}[nosep]
    \item Radical reform (shock therapy)
    \item EU anchor providing external $\Psi$ stability
    \item Digital leapfrog strategy
    \item \textbf{Result:} Accepted 5 years of $K \ll 1$ to reach high-$Q$ equilibrium
\end{itemize}

\textbf{Ukraine path:}
\begin{itemize}[nosep]
    \item Partial reform
    \item Oligarch capture
    \item No external anchor
    \item \textbf{Result:} Permanent valley—never achieved coherent configuration
\end{itemize}

\textbf{Insight:} Partial reform is worse than full reform or no reform.
The ``valley'' must be crossed quickly.

\end{tcolorbox}

\subsection{Level 6 — Global: ``Why Can't We Coordinate on Climate?''}

\begin{tcolorbox}[colback=orange!5!white, colframe=orange!75!black,
    title=\textbf{Worked Example L=6: Global Coordination Failure}]

\textbf{Situation:} All countries would benefit from limiting warming.
Yet coordination fails repeatedly.

\textbf{Analysis:}
\begin{itemize}[nosep]
    \item Each nation is internally coherent ($K_{\text{national}} \approx 1$)
    \item Global configuration is misfit: rhetoric of cooperation, structure of competition
    \item No global $\Psi$-setter (no world government)
\end{itemize}

\textbf{Why stuck:}
Free-rider problem. Short political time horizons. National identity conflicts with global identity.
$K_{\text{global}} \approx 0$ but stable because no actor can unilaterally change structure.

\textbf{Insight:} Climate failure is a \emph{coherent bad equilibrium} at global level.
Structure trumps intention. $K_{\text{local}}$ high, $K_{\text{global}}$ low.

\end{tcolorbox}

%==============================================================================
% PART III: SYNTHESIS
%==============================================================================

\subsection{Universal Pattern Across All Levels}

\begin{tcolorbox}[colback=blue!5!white, colframe=blue!75!black,
    title=\textbf{Integration: The Three Laws of Coherence}]

Across all six levels, the same patterns emerge:

\begin{enumerate}
    \item \textbf{Coherence $\neq$ Quality:} $K \approx 1$ does not imply high $Q$.
          Stable misery (``coherence trap'') is possible at every level.

    \item \textbf{Partial change is dangerous:} Moving strategy without structure
          (or vice versa) creates misfit. The ``valley of incoherence'' is costly.

    \item \textbf{Transformation requires commitment:} Successful change means
          accepting temporary $K < 1$ to reach a new, better $C^*$.
\end{enumerate}

\textbf{Implication for intervention:} Don't optimize within a bad equilibrium.
Either accept current $C^*$ or commit to full transformation through the valley.

\end{tcolorbox}

%==============================================================================
% BACK MATTER
%==============================================================================

\subsection{Summary}

\begin{tcolorbox}[colback=gray!5!white, colframe=gray!75!black,
    title=\textbf{Summary: Appendix EXM Key Points}]

\begin{enumerate}
    \item Six worked examples demonstrate EBF at each aggregation level
    \item $K$ (coherence) and $Q$ (quality) are independent dimensions
    \item Stable bad equilibria exist at all levels—the ``coherence trap''
    \item Partial transformation creates persistent misfit
    \item Successful change requires crossing the ``valley of incoherence''
    \item Same structural patterns appear from individual to global level
\end{enumerate}

\end{tcolorbox}

%------------------------------------------------------------------------------
% GLOSSARY SECTION
%------------------------------------------------------------------------------

\subsection{Key Terms}

\begin{tabular}{@{}p{3cm}p{9cm}@{}}
\textbf{Term} & \textbf{Definition} \\
\hline
Coherence ($K$) & Degree of alignment between strategy, structure, and context \\
Quality ($Q$) & Welfare level of a configuration \\
Coherence Trap & Stable bad equilibrium with high $K$ but low $Q$ \\
Valley of Incoherence & Temporary low-$K$ phase during transformation \\
Misfit & Configuration with $K \ll 1$ due to strategy-structure mismatch \\
\end{tabular}

\noindent\textit{See Appendix~G for complete definitions.}

%------------------------------------------------------------------------------
% CRITICAL FOUNDATIONS
%------------------------------------------------------------------------------

\subsection{Critical Foundations}

\begin{tcolorbox}[colback=red!5!white, colframe=red!75!black,
    title=\textbf{Critical Foundations for Appendix EXM}]

This appendix's worked examples depend on:

\begin{enumerate}
    \item \textbf{Appendix WHO (CORE-WHO):} Definition of aggregation levels $L$
    \item \textbf{Appendix HOW (CORE-HOW):} Complementarity structure creating $K$
    \item \textbf{Appendix PRF (FORMAL-PROOF):} Formal definitions of $K$, $Q$, $C^*$
    \item \textbf{Chapter 6:} Theory of reference structures
\end{enumerate}

\end{tcolorbox}

%------------------------------------------------------------------------------
% OPEN ISSUES
%------------------------------------------------------------------------------

\subsection{Open Issues}

\begin{enumerate}
    \item \textbf{Quantification:} How to measure $K$ and $Q$ empirically at each level?
    \item \textbf{Transition dynamics:} How long must the ``valley'' be endured?
    \item \textbf{Level interactions:} How does $K$ at one level affect $K$ at others?
    \item \textbf{Additional examples:} More domain-specific cases needed (health, education, finance)
\end{enumerate}

%------------------------------------------------------------------------------
% REFERENCES
%------------------------------------------------------------------------------

\subsection{References}

\noindent\textit{This appendix applies concepts developed in:}

\begin{itemize}[nosep]
    \item Chapter 6 (Reference Structure)
    \item Chapter 7 (Fit and Non-Concavity)
    \item Appendix PRF (Formal Proofs)
    \item Roberts (2004): Complementarity and organizational design
    \item Milgrom \& Roberts (1990): Modern manufacturing
\end{itemize}

\noindent\textit{For the master bibliography, see \texttt{appendices/bcm2\_master\_references.bib}.}

%%%%%%%%%%%%%%%%%%%%%%%%%%%%%%%%%%%%%%%%%%%%%%%%%%%%%%%%%%%%%%%%%%%%%%%%%%%%%%%


% Link to master bibliography
\nocite{bcm_master}


%%%%%%%%%%%%%%%%%%%%%%%%%%%%%%%%%%%%%%%%%%%%%%%%%%%%%%%%%%%%%%%%%%%%%%%%%%%%%%%

%%%%%%%%%%%%%%%%%%%%%%%%%%%%%%%%%%%%%%%%%%%%%%%%%%%%%%%%%%%%%%%%%%%%%%%%%%%%%%%
\section{Glossary of Symbols}
\label{sec:exm-glossary}
%%%%%%%%%%%%%%%%%%%%%%%%%%%%%%%%%%%%%%%%%%%%%%%%%%%%%%%%%%%%%%%%%%%%%%%%%%%%%%%

For comprehensive definitions, see Appendix G (Master Glossary).

\begin{table}[htbp]
\centering
\caption{Key Symbols in Appendix EXM}
\begin{tabular}{lll}
\toprule
\textbf{Symbol} & \textbf{Meaning} & \textbf{Reference} \\
\midrule
$\gamma$ & Complementarity parameter & Appendix B \\
$\Psi$ & Context vector & Appendix V \\
$U$ & Utility function & Chapter 3 \\
\bottomrule
\end{tabular}
\end{table}


\section{References}
\label{sec:exm-references}
%%%%%%%%%%%%%%%%%%%%%%%%%%%%%%%%%%%%%%%%%%%%%%%%%%%%%%%%%%%%%%%%%%%%%%%%%%%%%%%

See master bibliography (\texttt{bcm\_master.bib}) for complete citations.

\nocite{bcm_master}

